\documentclass{article}
\usepackage{propositions}
\usepackage{makeidx}\makeindex
\usepackage{hyperref}

%% align=runin puts the label into the body, so the item is already in
%% horizontal mode when whatever follows \item[...] is read.  That makes the
%% separating space vulnerable in a way it is not for the other alignments,
%% where the body has not begun and TeX drops stray spaces itself.
%%
%% \item[...] ends with \ignorespaces, which covers a space written directly
%% after it.  It does not cover a space after an intervening \label: that one
%% is left to \label's own \@esphack, which does \ignorespaces only if
%% \@savsk, the skip it found on entry, was positive.  Under hyperref the
%% anchor contributed by \refstepcounter used to sit between our space and the
%% \label, so \@savsk was zero, \@esphack declined, and the space after the
%% \label was typeset on top of ours: every run-in label in a hyperref
%% document carried a double space.  The separating space is now emitted after
%% the cross-referencing machinery, so it is the last thing on the list and
%% \@esphack does its job.
%%
%% This file loads hyperref for that reason.  test-runin.tex covers the rest
%% of the alignment's behaviour without it.

\begin{document}

\section{One space, whatever intervenes}

Every item below must show exactly one interword space after \texttt{Word},
matching the control line beneath the list.  A doubled space is the
regression this file exists to catch.

\noindent\rule{\textwidth}{0.3pt}
\begin{prop}
  \item[Word, align=runin] Nothing intervening.
  \item[Word, align=runin] \label{ri:1} A \verb|\label| after a space.
  \item[Word, align=runin]\label{ri:2} A \verb|\label| with no space before it.
  \item[Word, align=runin] \label{ri:3}
  A \verb|\label|, with the body on the next line.
  \item[Word, align=runin] \label{ri:4} \label{ri:5} Two labels.
  \item[Word, align=runin] \index{word} \label{ri:6} An \verb|\index| and a
  \verb|\label|.
  \item[Word, align=runin, label=ri:7] The \texttt{label} key rather than a
  \verb|\label|.
  \item[Word, align=runin]      Several spaces in the source.
\end{prop}
\noindent\rule{\textwidth}{0.3pt}

Word Control: this line's gap after \texttt{Word} is the one to match.

\section{The space still respects the space factor}

A label ending in a period or a colon must still take the wider space, since
the separator is an ordinary interword space and not a fixed skip.

\begin{prop}
  \item[{Word.}, align=runin] After a period.
  \item[{Word:}, align=runin] After a colon.
  \item[Word, align=runin] After a letter.
\end{prop}

The first two gaps must be visibly wider than the third.

\section{References are unaffected}

\ref{ri:1}, \ref{ri:3}, \ref{ri:7} --- ordinary numbered references; the
alignment governs placement only.

\end{document}
